% ============================================================
%  GUÍA 7: Introducción al Límite y Límites Algebraicos
%  Curso de Introducción a la Universidad — Precálculo y Cálculo I
%  Autor: Ing. Gabriel Astudillo
%  Clase cubierta: 7
% ============================================================

\documentclass[12pt, letterpaper]{article}

% ── Paquetes ─────────────────────────────────────────────────
\usepackage{mathwizards-guia}
\mwguia{Guía 7 — Límites Algebraicos}{Ing. Gabriel Astudillo $\cdot$ Curso Preuniversitario}

\begin{document}

% ── Portada / Título ─────────────────────────────────────────
\mwportada{Guía de Estudio 7}{Introducción al Límite y Límites Algebraicos}{Límite Intuitivo, Límites Laterales, Propiedades, Indeterminación $0/0$, Factorización y Racionalización}

\vspace{4pt}
\begin{cuadroinfo}
  \small
  \textbf{Presentación:} Esta guía construye la base conceptual del análisis matemático: la idea intuitiva de límite, los límites laterales y su teorema de existencia, las propiedades operativas y las técnicas para resolver indeterminaciones del tipo $0/0$ mediante factorización y racionalización. Incluye teoría, ejercicios resueltos paso a paso y propuestos con clave.
  \\[4pt]
  \textbf{Nivel de Dificultad:}\quad
  \facil{} Básico / Directo \quad
  \medio{} Intermedio \quad
  \dificil{} Desafiante \quad
  \muyDificil{} Muy Difícil / Reto
\end{cuadroinfo}

\vspace{8pt}

\section{Introducción Intuitiva al Límite y Límites Laterales}
% ============================================================

\subsection{¿Qué es un límite?}

\begin{cuadroteoria}
  \textbf{Idea intuitiva:} El límite de $f(x)$ cuando $x$ se acerca a $a$, denotado $\lim_{x \to a} f(x)$, es el valor al que $f(x)$ se aproxima conforme $x$ se acerca a $a$ (por ambos lados), sin necesidad de que $f(a)$ esté definida.

  \textbf{Límites laterales:}
  \begin{itemize}[leftmargin=*]
    \item $\lim_{x \to a^-} f(x)$: límite por la \textbf{izquierda} (valores $x < a$).
    \item $\lim_{x \to a^+} f(x)$: límite por la \textbf{derecha} (valores $x > a$).
  \end{itemize}

  \textbf{Existencia del límite:}
  $$\boxed{\lim_{x \to a} f(x) = L \iff \lim_{x \to a^-} f(x) = L \;\text{ y }\; \lim_{x \to a^+} f(x) = L}$$
\end{cuadroteoria}

\begin{cuadrosolucion}
  \textbf{Observación importante:} El límite $\lim_{x \to a} f(x)$ \textbf{no depende} del valor de $f(a)$. La función puede estar definida o no en $x = a$, y el límite existe igual si los valores se aproximan a $L$.
\end{cuadrosolucion}

\subsection{Ejercicios resueltos}

\textbf{Ejemplo R1.} Evalúa $\lim_{x \to 2} (x^2 + 1)$ mediante una tabla de valores.

\begin{cuadrosolucion}
  \textbf{Solución:} Construimos una tabla con valores cercanos a 2:
  \begin{center}
    \small
    \begin{tabular}{cccccccc}
      \toprule
      $x$            & 1{,}9  & 1{,}99   & 1{,}999 & $\to 2^-$ & $\to 2^+$ & 2{,}001 & 2{,}01   \\
      \midrule
      $f(x) = x^2+1$ & 4{,}61 & 4{,}9601 & 4{,}996 & $\to 5$   &           & 5{,}004 & 5{,}0401 \\
      \bottomrule
    \end{tabular}
  \end{center}
  Los valores se acercan a 5 por ambos lados. Por tanto: $\lim_{x \to 2} (x^2 + 1) = 5$.
\end{cuadrosolucion}
\textbf{Ejemplo R2.} Sea $f(x) = \dfrac{x^2 - 1}{x - 1}$ para $x \neq 1$. Evalúa $\lim_{x \to 1} f(x)$.

\begin{cuadrosolucion}
  \textbf{Solución:} Nota que $f(1)$ no está definida (división entre cero). Pero:
  $$\frac{x^2 - 1}{x - 1} = \frac{(x - 1)(x + 1)}{x - 1} = x + 1 \quad (x \neq 1).$$
  Entonces, cuando $x$ se acerca a 1, $f(x) = x + 1$ se acerca a 2. Por tanto:
  $$\lim_{x \to 1} \frac{x^2 - 1}{x - 1} = 2.$$
  Este es un ejemplo de \textbf{discontinuidad evitable}: la función tiene un hueco en $(1, 2)$, pero el límite existe.
\end{cuadrosolucion}

\textbf{Ejemplo R3.} Analiza la existencia del límite para $f(x) = |x|/x$ cuando $x \to 0$.

\begin{cuadrosolucion}
  \textbf{Solución:} Para $x > 0$: $\dfrac{|x|}{x} = 1$. Para $x < 0$: $\dfrac{|x|}{x} = -1$.
  \begin{itemize}
    \item $\lim_{x \to 0^+} \dfrac{|x|}{x} = 1$.
    \item $\lim_{x \to 0^-} \dfrac{|x|}{x} = -1$.
  \end{itemize}
  Como los límites laterales son distintos, $\lim_{x \to 0} \dfrac{|x|}{x}$ \textbf{no existe}.
\end{cuadrosolucion}

\textbf{Ejemplo R4.} Sea $f(x) = \begin{cases} x^2 + 1 & \text{si } x \ge 1 \\ 3x & \text{si } x < 1 \end{cases}$. ¿Existe $\lim_{x \to 1} f(x)$?

\begin{cuadrosolucion}
  \textbf{Solución:}
  \begin{itemize}
    \item Por la izquierda: $\lim_{x \to 1^-} 3x = 3$.
    \item Por la derecha: $\lim_{x \to 1^+} (x^2 + 1) = 2$.
  \end{itemize}
  Son distintos, por lo que el límite no existe.
\end{cuadrosolucion}

\subsection{Ejercicios propuestos}

\begin{enumerate}[series=ejercicios, leftmargin=*]
  \item \facil{} Evalúa mediante tabla de valores: $\lim_{x \to 3} (2x + 1)$.

  \item \facil{} Evalúa mediante tabla de valores: $\lim_{x \to 0} \cos x$.

  \item \facil{} Evalúa: $\lim_{x \to 2} \dfrac{x^2 - 4}{x - 2}$ (pista: factoriza).

  \item \facil{} Determina si el límite existe: $f(x) = \begin{cases} x + 2 & \text{si } x < 1 \\ 4 & \text{si } x = 1 \\ 2x + 1 & \text{si } x > 1 \end{cases}$, en $x = 1$.

  \item \medio{} Determina si existe $\lim_{x \to 0} \dfrac{1}{x^2}$.

  \item \medio{} Determina si existe $\lim_{x \to 2} \dfrac{2x - 4}{|x - 2|}$.

  \item \medio{} Sea $f(x) = \begin{cases} x^2 & \text{si } x < 0 \\ 1 & \text{si } x = 0 \\ x + 2 & \text{si } x > 0 \end{cases}$. Determina si existe $\lim_{x \to 0} f(x)$.

  \item \medio{} Determina si existe $\lim_{x \to 4} \dfrac{x - 4}{\sqrt{x} - 2}$ (pista: racionaliza).

  \item \dificil{} Dibuja una función que cumpla simultáneamente: $f(2) = 5$, $\lim_{x \to 2^-} f(x) = 3$, $\lim_{x \to 2^+} f(x) = 3$ y $\lim_{x \to 3} f(x) = \infty$.

  \item \dificil{} Encuentra una fórmula para una función que tenga un hueco en $x = -1$, salto en $x = 0$ y esté definida en todos los reales.

  \item \dificil{} Sea $f(x) = \begin{cases} 3x - a & \text{si } x < 2 \\ ax + 1 & \text{si } x \ge 2 \end{cases}$. Encuentra $a$ para que $\lim_{x \to 2} f(x)$ exista.

  \item \dificil{} Analiza la existencia de $\lim_{x \to 1} \dfrac{|x - 1|}{x - 1}$ y justifica tu respuesta.
\end{enumerate}

\newpage

% ============================================================

\section{Cálculo de Límites Algebraicos}
% ============================================================

\subsection{Propiedades de los límites}

\begin{cuadroteoria}
  Si $\lim_{x \to a} f(x) = L$ y $\lim_{x \to a} g(x) = M$, entonces:
  \begin{align*}
    \lim_{x \to a} [f(x) + g(x)]                  & = L + M                                                               \\
    \lim_{x \to a} [f(x) - g(x)]                  & = L - M                                                               \\
    \lim_{x \to a} [f(x) \cdot g(x)]              & = L \cdot M                                                           \\
    \lim_{x \to a} \left[\frac{f(x)}{g(x)}\right] & = \frac{L}{M} \quad (M \neq 0)                                        \\
    \lim_{x \to a} [c \cdot f(x)]                 & = c \cdot L                                                           \\
    \lim_{x \to a} [f(x)]^n                       & = L^n                                                                 \\
    \lim_{x \to a} \sqrt[n]{f(x)}                 & = \sqrt[n]{L} \quad (\text{si } L \geq 0 \text{ para } n \text{ par})
  \end{align*}

  \textbf{Límite de una función polinomial:}
  $$\lim_{x \to a} P(x) = P(a)$$

  \textbf{Sustitución directa:} para funciones polinomiales y racionales (cuando el denominador no se anula), el límite se evalúa sustituyendo $x = a$.
\end{cuadroteoria}

\begin{cuadroadvertencia}
  \textbf{Forma indeterminada $0/0$:} Si al sustituir obtienes $\frac{0}{0}$, necesitas \textbf{simplificar} antes de evaluar. Técnicas:
  \begin{enumerate}[leftmargin=*]
    \item \textbf{Factorizar} el numerador y/o denominador, y cancelar el factor común.
    \item \textbf{Racionalizar} (multiplicar por el conjugado) cuando hay raíces.
    \item En límites de cocientes de polinomios, el factor que causa el problema siempre se cancela o se llega a un límite infinito.
  \end{enumerate}
\end{cuadroadvertencia}

\subsection{Ejercicios resueltos}

\textbf{Ejemplo R1.} Calcula $\lim_{x \to 3} (2x^2 - 5x + 1)$.

\begin{cuadrosolucion}
  \textbf{Solución:} Sustitución directa:
  $$2(3)^2 - 5(3) + 1 = 18 - 15 + 1 = 4.$$
\end{cuadrosolucion}

\textbf{Ejemplo R2.} Calcula $\lim_{x \to 2} \dfrac{x^2 - 5x + 6}{x - 2}$.

\begin{cuadrosolucion}
  \textbf{Solución:} Sustituyendo $x = 2$ llegamos a $0/0$. Factorizamos:
  $$\frac{x^2 - 5x + 6}{x - 2} = \frac{(x - 2)(x - 3)}{x - 2} = x - 3 \quad (x \neq 2).$$
  Entonces: $\lim_{x \to 2} \dfrac{x^2 - 5x + 6}{x - 2} = \lim_{x \to 2} (x - 3) = -1$.
\end{cuadrosolucion}

\textbf{Ejemplo R3.} Calcula $\lim_{x \to 4} \dfrac{\sqrt{x} - 2}{x - 4}$.

\begin{cuadrosolucion}
  \textbf{Solución:} Sustituyendo da $0/0$. Racionalizamos el numerador multiplicando por el conjugado $\sqrt{x} + 2$:
  $$\frac{(\sqrt{x} - 2)(\sqrt{x} + 2)}{(x - 4)(\sqrt{x} + 2)} = \frac{x - 4}{(x - 4)(\sqrt{x} + 2)} = \frac{1}{\sqrt{x} + 2}.$$
  Entonces:
  $$\lim_{x \to 4} \dfrac{\sqrt{x} - 2}{x - 4} = \lim_{x \to 4} \dfrac{1}{\sqrt{x} + 2} = \dfrac{1}{\sqrt{4} + 2} = \dfrac{1}{4}.$$
\end{cuadrosolucion}

\textbf{Ejemplo R4.} Calcula $\lim_{x \to 1} \dfrac{x^3 - 1}{x - 1}$.

\begin{cuadrosolucion}
  \textbf{Solución:} Factorizamos la diferencia de cubos $x^3 - 1 = (x - 1)(x^2 + x + 1)$:
  $$\frac{x^3 - 1}{x - 1} = x^2 + x + 1 \quad (x \neq 1).$$
  $$\lim_{x \to 1} \dfrac{x^3 - 1}{x - 1} = \lim_{x \to 1} (x^2 + x + 1) = 1 + 1 + 1 = 3.$$
\end{cuadrosolucion}

\textbf{Ejemplo R5.} Calcula $\lim_{x \to 0} \dfrac{\sqrt{1 + x} - 1}{x}$.

\begin{cuadrosolucion}
  \textbf{Solución:} Racionalizamos:
  $$\frac{(\sqrt{1 + x} - 1)(\sqrt{1 + x} + 1)}{x(\sqrt{1 + x} + 1)} = \frac{1 + x - 1}{x(\sqrt{1 + x} + 1)} = \frac{x}{x(\sqrt{1 + x} + 1)} = \frac{1}{\sqrt{1 + x} + 1}.$$
  $$\lim_{x \to 0} \dfrac{\sqrt{1 + x} - 1}{x} = \dfrac{1}{\sqrt{1} + 1} = \dfrac{1}{2}.$$
  Este límite aparecerá de nuevo en la definición de derivada.
\end{cuadrosolucion}

\newpage

\subsection{Ejercicios propuestos}

\begin{enumerate}[resume=ejercicios, leftmargin=*]
  \item \facil{} Calcula: $\lim_{x \to 1} (x^3 - 2x^2 + x)$.

  \item \facil{} Calcula: $\lim_{x \to -2} \dfrac{x^2 + x - 2}{x + 2}$.

  \item \facil{} Calcula: $\lim_{x \to 5} \dfrac{x^2 - 25}{x - 5}$.

  \item \facil{} Calcula: $\lim_{x \to 3} \dfrac{2x - 6}{x^2 - 9}$.

  \item \medio{} Calcula: $\lim_{x \to 9} \dfrac{\sqrt{x} - 3}{x - 9}$.

  \item \medio{} Calcula: $\lim_{x \to 0} \dfrac{\sqrt{4 + x} - 2}{x}$.

  \item \medio{} Calcula: $\lim_{x \to 2} \dfrac{x^3 - 8}{x - 2}$.

  \item \medio{} Calcula: $\lim_{x \to 1} \dfrac{x^2 + 2x - 3}{x^2 - 1}$.

  \item \dificil{} Calcula: $\lim_{x \to 0} \dfrac{1 - \sqrt{1 - x}}{x}$.

  \item \dificil{} Calcula: $\lim_{x \to 1} \dfrac{\sqrt{x} - 1}{x - 1}$ (pista: haz $t = \sqrt{x}$, con $t \to 1$).

  \item \dificil{} Calcula: $\lim_{h \to 0} \dfrac{(2 + h)^3 - 8}{h}$.

  \item \dificil{} Calcula: $\lim_{x \to 0} \dfrac{3x - \sin(3x)}{x}$ (pista: usa el límite fundamental que verás en la siguiente sección).
\end{enumerate}

\newpage

% ============================================================

% ──────────────────────────────────────────────────────────
%  NUEVOS EJERCICIOS PROPUESTOS (26) — INSERCIÓN MANUAL
% ──────────────────────────────────────────────────────────
\subsection{Ejercicios propuestos: Límites Intuitivos y Laterales}

\begin{enumerate}[resume=ejercicios, leftmargin=*]
  \item \facil{} Calcula el límite por sustitución directa: $\displaystyle\lim_{x \to 2} (3x + 5)$.

  \item \facil{} Calcula el límite simplificando: $\displaystyle\lim_{x \to 1} \frac{x^2 - 1}{x - 1}$.

  \item \medio{} Calcula el límite factorizando: $\displaystyle\lim_{x \to 3} \frac{x^2 - 9}{x - 3}$.

  \item \medio{} Calcula el límite racionalizando: $\displaystyle\lim_{x \to 0} \frac{\sqrt{x + 1} - 1}{x}$.

  \item \medio{} Determina el límite lateral: $\displaystyle\lim_{x \to 2^-} \frac{1}{x - 2}$.

  \item \dificil{} Calcula el límite racionalizando: $\displaystyle\lim_{x \to 4} \frac{\sqrt{x} - 2}{x - 4}$.

  \item \dificil{} Calcula el límite factorizando la diferencia de cubos: $\displaystyle\lim_{x \to 1} \frac{x^3 - 1}{x - 1}$.

  \item \dificil{} Determina si existe $\displaystyle\lim_{x \to 0} \frac{|x|}{x}$ analizando los límites laterales.

  \item \muyDificil{} Calcula el límite con raíz en el numerador:
        $$\lim_{x \to 9} \frac{3 - \sqrt{x}}{9 - x}$$
\end{enumerate}

\subsection{Ejercicios propuestos: Límites Algebraicos Avanzados}

\begin{enumerate}[resume=ejercicios, leftmargin=*]
  \item \facil{} Calcula el límite simplificando: $\displaystyle\lim_{x \to 0} \frac{2x^2 + 3x}{x}$.

  \item \facil{} Calcula el límite por sustitución: $\displaystyle\lim_{x \to -1} (x^3 + 2x)$.

  \item \medio{} Calcula el límite factorizando: $\displaystyle\lim_{x \to 5} \frac{x^2 - 25}{x - 5}$.

  \item \medio{} Calcula el límite factorizando el trinomio: $\displaystyle\lim_{x \to 2} \frac{x^2 - 3x + 2}{x - 2}$.

  \item \medio{} Calcula el límite racionalizando: $\displaystyle\lim_{x \to 0} \frac{\sqrt{4 + x} - 2}{x}$.

  \item \dificil{} Calcula el límite factorizando el polinomio de grado 3:
        $$\lim_{x \to 1} \frac{x^3 - x^2 - x + 1}{x - 1}$$

  \item \dificil{} Calcula el límite racionalizando el denominador:
        $$\lim_{x \to 2} \frac{x - 2}{\sqrt{x + 2} - 2}$$

  \item \dificil{} Calcula el cociente incremental (derivada de $x^2$ en $x = 2$):
        $$\lim_{h \to 0} \frac{(2 + h)^2 - 4}{h}$$

  \item \muyDificil{} Calcula el límite con raíz cúbica:
        $$\lim_{x \to 1} \frac{\sqrt[3]{x} - 1}{x - 1}$$
\end{enumerate}

\subsection{Ejercicios propuestos: Retos Integradores}

\begin{enumerate}[resume=ejercicios, leftmargin=*]
  \item \facil{} Calcula el límite simplificando: $\displaystyle\lim_{x \to 3} \frac{x - 3}{x^2 - 9}$.

  \item \medio{} Calcula el límite expandiendo el cubo: $\displaystyle\lim_{x \to 0} \frac{(x + 1)^3 - 1}{x}$.

  \item \medio{} Calcula el límite con raíz: $\displaystyle\lim_{x \to 0} \frac{\sqrt{x^2 + 4} - 2}{x^2}$.

  \item \medio{} Calcula el límite lateral usando el valor absoluto:
        $$\lim_{x \to 0^-} \frac{|x| - x}{x}$$

  \item \dificil{} Calcula el límite con fracción algebraica:
        $$\lim_{x \to 1} \frac{\frac{1}{x} - 1}{x - 1}$$

  \item \dificil{} Calcula el límite racionalizando:
        $$\lim_{x \to 0} \frac{1 - \sqrt{1 - x}}{x}$$

  \item \muyDificil{} Calcula el límite con parámetro $a$ (constante real):
        $$\lim_{x \to a} \frac{x^2 - a^2}{x - a}$$

  \item \muyDificil{} Calcula el cociente incremental general (derivada de $\sqrt{x}$):
        $$\lim_{h \to 0} \frac{\sqrt{x + h} - \sqrt{x}}{h}$$
\end{enumerate}
\newpage
% ============================================================
\ifmwclaves
  \section{Clave de Respuestas}
  % ============================================================

  \subsection*{Sección 1: Introducción (Ejercicios 1--12)}

  \begin{enumerate}[series=respuestas, leftmargin=*, label=\textbf{\arabic*.}]
    \item $\lim_{x \to 3} (2x + 1) = 7$.

    \item $\lim_{x \to 0} \cos x = 1$.

    \item $\dfrac{x^2 - 4}{x - 2} = x + 2$, entonces el límite es $4$.

    \item Por izquierda: $3$; por derecha (y $f(1)$): $3$. Los límites laterales coinciden: el límite existe y vale 3.

    \item $x^2 \to 0^+$ siempre (es cuadrado), así que $\dfrac{1}{x^2} \to +\infty$. El límite es $+\infty$; no existe como límite finito.

    \item Por derecha: $\dfrac{2(x-2)}{x-2} = 2 \to 2$. Por izquierda: $\dfrac{2(x-2)}{-(x-2)} = -2 \to -2$. No existe (laterales distintos).

    \item Por izquierda: $\lim_{x \to 0^-} x^2 = 0$. Por derecha: $\lim_{x \to 0^+} (x + 2) = 2$. Laterales distintos, no existe.

    \item Racionalizamos: $\dfrac{(x - 4)(\sqrt{x} + 2)}{(\sqrt{x} - 2)(\sqrt{x} + 2)} = \dfrac{(x-4)(\sqrt{x}+2)}{x - 4} = \sqrt{x} + 2$. Límite: $\sqrt{4} + 2 = 4$.

    \item (Gráfico, respuesta abierta). Una función con hueco en $(2,3)$, salto en algún punto antes de 3, y asíntota vertical en $x = 3$.

    \item Ejemplo: $f(x) = \begin{cases} x + 2 & \text{si } x \neq -1, 0 \\ 1 & \text{si } x = 0 \end{cases}$ (hueco en $x = -1$: no definida; salto en $x = 0$). Otras respuestas son válidas.

    \item $\lim_{x \to 2^-} (3x - a) = 6 - a$; $\lim_{x \to 2^+} (ax + 1) = 2a + 1$. Igualando: $6 - a = 2a + 1 \Rightarrow 3a = 5 \Rightarrow a = \dfrac{5}{3}$.

    \item Para $x > 1$: $\dfrac{x-1}{x-1} = 1$; para $x < 1$: $\dfrac{-(x-1)}{x-1} = -1$. Laterales distintos, el límite no existe.
  \end{enumerate}


  \subsection*{Sección 2: Límites Algebraicos (Ejercicios 13--24)}

  \begin{enumerate}[resume=respuestas, leftmargin=*, label=\textbf{\arabic*.}]
    \item $1 - 2 + 1 = 0$.

    \item $\dfrac{x^2 + x - 2}{x + 2} = \dfrac{(x + 2)(x - 1)}{x + 2} = x - 1$. En $x = -2$: $-3$.

    \item $x + 5$; límite: $10$.

    \item $\dfrac{2(x - 3)}{(x - 3)(x + 3)} = \dfrac{2}{x + 3}$; en $x = 3$: $\dfrac{2}{6} = \dfrac{1}{3}$.

    \item $\dfrac{1}{\sqrt{x} + 3}$; en $x = 9$: $\dfrac{1}{6}$.

    \item $\dfrac{1}{\sqrt{4 + x} + 2}$; en $x = 0$: $\dfrac{1}{4}$.

    \item $x^2 + 2x + 4$; en $x = 2$: $12$.

    \item $\dfrac{(x - 1)(x + 3)}{(x - 1)(x + 1)} = \dfrac{x + 3}{x + 1}$; en $x = 1$: $\dfrac{4}{2} = 2$.

    \item Racionalizamos: $\dfrac{x}{(\sqrt{1 - x} + 1)(1 + \sqrt{1-x})} \dots$ Al final: $\lim_{x \to 0} \dfrac{1 - \sqrt{1 - x}}{x} = \dfrac{1}{2}$.

    \item Con $t = \sqrt{x}$: $x = t^2$, cuando $x \to 1$, $t \to 1$: $\dfrac{t - 1}{t^2 - 1} = \dfrac{1}{t + 1}$. Límite: $\dfrac{1}{2}$.

    \item Expandimos: $(2 + h)^3 = 8 + 12h + 6h^2 + h^3$. Restando 8 y dividiendo entre $h$: $12 + 6h + h^2$. Límite: $12$.

    \item $\dfrac{3x - \sin(3x)}{x} = 3 - \dfrac{\sin(3x)}{x} = 3 - 3\left(\dfrac{\sin(3x)}{3x}\right) \to 3 - 3 = 0$.
  \end{enumerate}


  % ──────────────────────────────────────────────────────────
  %  NUEVAS RESPUESTAS (26) — INSERCIÓN MANUAL
  % ──────────────────────────────────────────────────────────
  \subsection*{Sección 3: Límites Intuitivos y Laterales (Ejercicios 25--33)}
  \begin{enumerate}[resume=respuestas, leftmargin=*, label=\textbf{\arabic*.}]
    \item $11$
    \item $\displaystyle\lim_{x \to 1} (x + 1) = 2$
    \item $\displaystyle\lim_{x \to 3} (x + 3) = 6$
    \item $\displaystyle\lim_{x \to 0} \frac{1}{\sqrt{x + 1} + 1} = \frac{1}{2}$
    \item $-\infty$ (el denominador tiende a 0 por valores negativos)
    \item $\displaystyle\lim_{x \to 4} \frac{1}{\sqrt{x} + 2} = \frac{1}{4}$
    \item $\displaystyle\lim_{x \to 1} (x^2 + x + 1) = 3$
    \item No existe: el límite por la izquierda es $-1$ y por la derecha es $+1$.
    \item $\displaystyle\lim_{x \to 9} \frac{1}{3 + \sqrt{x}} = \frac{1}{6}$
  \end{enumerate}

  \subsection*{Sección 4: Límites Algebraicos Avanzados (Ejercicios 34--42)}
  \begin{enumerate}[resume=respuestas, leftmargin=*, label=\textbf{\arabic*.}]
    \item $\displaystyle\lim_{x \to 0} (2x + 3) = 3$
    \item $(-1)^3 + 2(-1) = -3$
    \item $\displaystyle\lim_{x \to 5} (x + 5) = 10$
    \item $\displaystyle\lim_{x \to 2} (x - 1) = 1$
    \item $\displaystyle\lim_{x \to 0} \frac{1}{\sqrt{4 + x} + 2} = \frac{1}{4}$
    \item Factorizando: $(x - 1)^2 (x + 1)$; límite $= 0$
    \item $\displaystyle\lim_{x \to 2} \big(\sqrt{x + 2} + 2\big) = 4$
    \item $\displaystyle\lim_{h \to 0} (4 + h) = 4$
    \item Con $a^3 - b^3$: $\displaystyle\lim_{x \to 1} \frac{1}{\sqrt[3]{x^2} + \sqrt[3]{x} + 1} = \frac{1}{3}$
  \end{enumerate}

  \subsection*{Sección 5: Retos Integradores (Ejercicios 43--50)}
  \begin{enumerate}[resume=respuestas, leftmargin=*, label=\textbf{\arabic*.}]
    \item $\displaystyle\lim_{x \to 3} \frac{1}{x + 3} = \frac{1}{6}$
    \item $\displaystyle\lim_{x \to 0} (3 + 3x + x^2) = 3$
    \item $\displaystyle\lim_{x \to 0} \frac{1}{\sqrt{x^2 + 4} + 2} = \frac{1}{4}$
    \item Para $x < 0$: $\dfrac{-x - x}{x} = -2$
    \item $\displaystyle\lim_{x \to 1} \left(-\frac{1}{x}\right) = -1$
    \item $\displaystyle\lim_{x \to 0} \frac{1}{1 + \sqrt{1 - x}} = \frac{1}{2}$
    \item $\displaystyle\lim_{x \to a} (x + a) = 2a$
    \item $\displaystyle\lim_{h \to 0} \frac{1}{\sqrt{x + h} + \sqrt{x}} = \frac{1}{2\sqrt{x}}$
  \end{enumerate}
\fi
\end{document}
